\section{Discussion}\label{sec:discussion}
% Interpretation of results
% Comparison with existing literature
% Limitations of this review
% Future research directions

\subsection{Interpretation of the Main Findings}\label{subsec:interpretation-of-the-main-findings}

The successful application of Machine Learning (ML) to the monitoring of Equatorial Plasma Bubbles (EPBs) signifies a paradigm shift from labor-intensive, subjective manual observation to automated, highly scalable analytical frameworks \cite{zhong2025automatic,okoh2025bootstrapping}. The findings demonstrate that deep learning and advanced feature extraction algorithms do not merely replicate human classification but can identify subtle spatial gradients, textures, and intensity variations that are often visually ambiguous to the human eye \cite{okoh2025bootstrapping}. Moreover, the automated extraction of morphological parameters has proven robust enough to generate solar-cycle-scale statistical profiles that were previously impossible to compile efficiently \cite{zhong2025statistical}.

However, the overarching limitations reported across these studies consistently point toward a strong susceptibility to environmental contamination and the systemic challenge of handling heavily imbalanced datasets. These structural trends establish clear performance trade-offs and operational boundary conditions, which are synthesized in detail in the subsequent sections.

\subsection{Comparison Between Machine Learning Techniques}\label{subsec:comparison-between-machine-learning-techniques}

A critical analysis of the deployed architectures reveals a distinct trade-off between computational complexity and operational efficiency. Deep Convolutional Neural Networks (DCNNs) such as ResNet and YOLO are highly effective at capturing the non-rigid, branching morphologies of EPBs and performing pixel-level semantic segmentation. \cite{zhong2025automatic,srisamoodkham2022detecting} However, increasing the depth of these networks often yields only marginal improvements in accuracy while significantly increasing processing latency, making them computationally expensive \cite{thanakulketsarat2023classification}.

Conversely, hybrid models integrating dimensionality reduction with traditional classifiers---such as Two-Dimensional Principal Component Analysis (2DPCA) or Singular Value Decomposition (SVD) combined with Support Vector Machines (SVM) or Random Forests (RF)---offer a highly efficient alternative. By directly processing the 2D image matrix and eliminating redundant features, these techniques dramatically reduce trained model size (by up to 70\%) and training time without compromising detection accuracy \cite{yacoub2025automatic,thanakulketsarat2023classification,siddiqui2025enhanced}. Additionally, the use of ensemble techniques, such as bootstrapping CNNs with probability averaging or mode voting, mitigates the overconfidence of individual sub-models and enhances robustness against ambiguous cases \cite{okoh2025bootstrapping}. Finally, the integration of Explainable AI (XAI) frameworks, specifically SHAP and LIME, resolves the ``black-box'' dilemma of deep learning by confirming that automated classifiers rely on physically meaningful features rather than background artifacts \cite{siddiqui2025enhanced}.

\subsection{Impact of Different Data Types}\label{subsec:impact-of-different-data-types}

The efficacy of ML models is intrinsically linked to the data modality utilized. Optical All-Sky Imagers (ASI) remain the primary data source due to their wide field of view and operational ability to capture fine-scale 2D morphologies and bifurcation structures \cite{githio2024machine,zhong2025automatic}. In contrast, Global Navigation Satellite System (GNSS) data, utilized through Total Electron Content (TEC) maps and ROTI indices, offers weather-independent, continuous global coverage \cite{githio2024machine}. However, the spatial resolution of TEC maps is generally inferior to ASI for observing the highly localized, fine structures of EPBs. Very High Frequency (VHF) radar provides exceptionally high spatial and temporal resolution for observing bottom-side F-layer boundary irregularities, but it is constrained by a localized geographical footprint \cite{thanakulketsarat2023classification}. The complementary nature of these data types suggests that relying on a single instrument limits the scope of monitoring systems \cite{githio2024machine}.

\subsection{Recurring Limitations in the Studies}\label{subsec:recurring-limitations-in-the-studies}

Despite high reported accuracies, several recurring limitations persist across the literature. Foremost is the \textbf{severe susceptibility to environmental contamination}. While ASI provides superior resolution, its inherent reliance on optical visibility means that tropospheric weather conditions---such as heavy cloud cover, lunar illumination, or artificial light pollution---can obscure EPB signatures. This drastically reduces the reliability of optical-based models and frequently leads to misclassifications \cite{yacoub2025automatic,siddiqui2025enhanced}.

Another critical limitation is the \textbf{imbalance in training datasets}. EPB occurrence is highly correlated with solar activity, resulting in datasets skewed heavily toward high solar activity years \cite{zhong2025automatic,okoh2025bootstrapping}. Consequently, models tend to become biased toward the majority class, performing optimally during solar maxima or geomagnetically quiet periods, but struggling during solar minima or severe geomagnetic storms where sample sizes are much smaller \cite{zhong2025automatic}.

Furthermore, because these are predominantly supervised learning models, they rely entirely on \textbf{manual ground-truth annotations}. The subjective nature of manually distinguishing faint EPB signatures from background noise or clouds establishes an artificial ceiling on model accuracy, as human error is inadvertently embedded into the training data \cite{okoh2025bootstrapping}.

\subsection{Challenges in the Field}\label{subsec:challenges-in-the-field}

A major challenge lies in the \textbf{generalizability of the models}. Models trained on data from a specific geographical sector (e.g., Brazil, China, or Thailand) or a specific imager configuration may experience degraded performance when applied to different locations due to varying camera characteristics, viewing geometries, and regional EPB morphologies \cite{siddiqui2025enhanced}. Additionally, balancing real-time operational processing speeds with the computational demands of deep learning algorithms remains a hurdle for deploying these systems on low-cost, edge-computing devices \cite{thanakulketsarat2023classification,yacoub2025automatic}.

\subsection{Gaps in the Literature}\label{subsec:gaps-in-the-literature}

Current research is heavily skewed toward reactive classification and detection. There is a \textbf{stark lack of predictive models} capable of forecasting the initiation, growth rate, and future trajectory of EPBs before they fully manifest \cite{siddiqui2025enhanced}. Furthermore, while researchers often use multiple instruments (e.g., ASI and GNSS) to validate findings independently, there is an \textbf{absence of true multimodal deep learning frameworks} that ingest and fuse optical, radar, and TEC data simultaneously within a single neural network architecture to overcome individual instrumental blind spots \cite{siddiqui2025enhanced}.

\subsection{Opportunities for Future Work}\label{subsec:opportunities-for-future-work}

Future research should prioritize the development of \textbf{multi-instrument data fusion models}, integrating ASI, GNSS-TEC, and satellite in-situ measurements to create robust systems that maintain high accuracy regardless of tropospheric weather conditions \cite{siddiqui2025enhanced}. Additionally, there is a critical need to transition from detection to \textbf{spatiotemporal predictive modeling}. By applying sequence-based networks (e.g., LSTMs or Transformers) to time-series airglow imagery and geophysical indices, models could proactively forecast EPB zonal drift and morphological evolution \cite{siddiqui2025enhanced}. Finally, XAI should be leveraged not just for model validation, but for \textbf{scientific discovery}, using interpretability tools to uncover previously unnoticed physical relationships and environmental drivers governing ionospheric dynamics \cite{siddiqui2025enhanced}.
